% Transcription — fonds Alexandre Grothendieck, shelfmark 11, batch 2 (pages 21-40).
% Established from the facsimile archives/batches/11/batch-02.pdf (a mirror of
% https://grothendieck.umontpellier.fr/11.pdf), per the skill
% transcribe-grothendieck. PDF page 1 of the batch file is archivists' page 21
% (the batch file carries no cover sheet; checked against the pencilled numbers).
%
% Pass: Opus 5 (claude-opus-5), 2026-09-14 — first pass, unchecked against the pages by a human.
%
% Nine of the twenty pages are his. Pages 21, 23, 25, 27, 29, 31, 34, 36, 38
% and 40 are pages of an unrelated English typescript (several scanned upside
% down) and page 32 is a blank verso through which page 33 shows; all eleven are
% skipped. Every leaf of his in the batch is fast drafting: the formulas and
% the underlined statements carry, much of the connective prose does not.
%
%   22  « Sous-anneaux bivariants sur V », taking up the stable sub-functors
%       A' ⊂ A of page 20: intersections, the stable sub-functor A_sp
%       generated, the NB that A_sp(X) contains the f_*(1_Y) and in many cases
%       (smooth projective schemes; char. 0 granted resolution of
%       singularities) is exactly their set; applications « du type A' »,
%       x_1,…,x_n ↦ u(x_1 ⊠ … ⊠ x_n) with u ∈ A'(X_1×…×X_n×X), and their
%       composites.
%   24  composites of maps of type A' are of type A'; « Application »: the
%       smallest stable sub-functor B ⊂ A containing A' and given E_X ⊂ A(X)
%       is the set of the φ(x_1,…,x_n), φ A'-admissible, x_i ∈ E(X_i).
%   26  a sign computation for the symmetry s_{A,B} against f×g, cancelled by
%       two diagonal strokes; then section « A »: the category V° of Poincaré
%       algebras (graded, finite type, degrees ≥ 0, a form κ_A : A → k of degree
%       d°(A) with (x,y) ↦ κ_A(xy) strictly non-degenerate; morphisms plain
%       graded algebra homs); the contravariant functor φ : V → Modgrad(k),
%       φ(f*) = f_* defined by κ_A(f*(b)a) = κ_B(b f_*(a)), projection
%       formula and κ_B f_* = κ_A.
%   28  hence a functor V → graded bivariant algebras; 3) k initial in V°,
%       A ⊗ B the coproduct (product in V), κ_{A⊗B}(x⊗y) = κ_A(x)κ_B(y)(−1)^{Bx};
%       the formula (f×g)_* = f_* ⊗ g_* « est fausse » unless f, g are of even
%       degree; the sign bookkeeping of four conventions, ending on the boxed
%       (f×g)_* = (f_* ⊗^{ord} g_*)(id_A ⊗ [(−1)^B]^{deg f}) and « Il n'y a
%       aucun signe dès que f est pair ».
%   30  for f of even degree (f×g)_* = f_* ⊗^{ord} g_*; the symmetry
%       s(x⊗y) = (−1)^{xy} y⊗x is compatible; the base-change formulas for the
%       projections X×S → X, checked « O.K. »; the composition X → Y → Z through
%       graphs Γ_f, Γ_{gf}, λ = Γ_f × id_Z, λ' = id_X × Γ_g; closing « Question »:
%       when does the exchange formula hold — is it enough that the fibre
%       product exist in V with the right dimensions?
%   33  « Catégorie des algèbres bivariantes (sur un anneau k) »: pairs
%       f = (f_*, f*), f* an algebra hom, f_* B-linear (projection formula), and
%       κ_B f_* = κ_A; marginal « Remarque » on isomorphisms; for A with
%       φ_A : A → Ǎ, ⟨y, φ_A(x)⟩ = κ(yx), and φ_B injective, f_* is determined by
%       f* through (***) κ(y f_*(x)) = κ(f*(y)x).
%   35  with κ_A factoring through A^{d(A)}, f_* raises degrees by d' − d; the
%       working category: graded anticommutative bivariant algebras with a
%       degree d, f_* raising degrees by d(B) − d(A), κ_A through A^{d(A)};
%       « Il n'est pas question de se borner aux algèbres de Poincaré ».
%   37  the graded tensor product is a functor on C but not a product in the
%       bivariant category (no diagonal), except for Poincaré algebras;
%       simplicial bivariant algebras (« pas ½ simpliciale ! »), contravariant
%       on finite non-empty sets; « spéciales » when A(I')⊗A(I'') ⥲ A(I'⊔I'');
%       with A_0 Poincaré, an equivalence with Poincaré algebras.
%   39  « Proposition »: for sub-algebras A_n ⊂ H*(X^n) stable under the
%       symmetric group, equivalence of (i) closure under correspondences u(x),
%       ⊠, graphs of surjections [1,n+1] → [1,n] (and injections), (ii) closure
%       under inverse and direct images by projections and injections, and
%       (iii) the A_{n+m}, read as sets of homs H*(X^n) → H*(X^m), closed under
%       composition, identities, ⊠ and simplicial maps.
%
% Notation fixed for the batch (kept in line with batches 1 and 3 of the
% folder): the category of the base \mathcal{V}, its opposite \mathcal{V}^{\circ}
% for Poincaré algebras; the external product \boxtimes; the form κ_A (his k_A,
% distinguished from the ground ring k as batch 3 does); the unit written U_Y on
% page 22 kept as $U_{Y}$; graphs Γ_f; ⊗^{ord} for his « ⊗ » surmounted by
% « ord ».
%
% Batch boundaries. Page 20 (batch 1) closes its proof with « Stabilité par ⊠
% claire » and page 22 opens on a new heading, so no sentence crosses the
% 20/21 boundary — but page 22 takes up the stable sub-functors A' ⊂ A of page
% 20, and says so by its notation. Page 39 ends complete on (iii) and page 40 is
% typescript; nothing visible crosses 40/41.

\documentclass[11pt,a4paper]{article}
% Le préambule commun à toutes les transcriptions du fonds Grothendieck.
%
% Il tient en une idée : **l'incertitude se compose, elle ne se résout pas.**
% Une transcription qui lisse un mot douteux a détruit la seule chose pour
% laquelle on la faisait — un lecteur doit pouvoir savoir, à chaque endroit,
% ce qui était sur le feuillet et ce qui a été ajouté. Les six macros qui
% suivent sont donc l'appareil critique, et non de la décoration.
%
% Les mêmes commandes sont reconnues par `scripts/render.mjs`, qui produit la
% vue de lecture du site. Une macro ajoutée ici doit l'être aussi là-bas,
% faute de quoi le rendu échouera bruyamment — ce qui est voulu : un rendu
% silencieusement faux est pire qu'un rendu absent.

\ProvidesPackage{grothendieck}[2026/08/08 Transcriptions du fonds Grothendieck]

\usepackage{amsmath,amssymb,amsthm}
\usepackage{mathrsfs}
\usepackage{stmaryrd}
% Les diagrammes commutatifs sont partout dans ce fonds : c'est la langue
% même de la géométrie algébrique de Grothendieck, et une transcription qui
% les rendrait en prose ne transcrirait rien.
\usepackage{tikz-cd}
% Français par défaut — c'est la langue des feuillets — mais l'anglais est
% chargé aussi : la traduction et le résumé sont des documents anglais, et la
% typographie française (espaces avant les deux-points) y serait une faute.
% Ces documents basculent par \selectlanguage{english} en tête de corps.
\usepackage[english,french]{babel}
\usepackage{fontspec}
\usepackage{geometry}
\geometry{a4paper,margin=25mm,marginparwidth=28mm}
\usepackage{marginnote}
\usepackage{xcolor}
% ulem, not soul. `soul`'s \st and \ul are built on a character-by-character
% scan that XeLaTeX's font machinery breaks: the document fails with an
% undefined control sequence rather than degrading. ulem does the same two jobs
% and is Unicode-engine safe. `normalem` stops it hijacking \emph, which the
% transcriptions use for Grothendieck's own emphasis.
\usepackage[normalem]{ulem}
\usepackage{hyperref}
\hypersetup{colorlinks=true,linkcolor=black,urlcolor=black,pdfborder={0 0 0}}

% --- Métadonnées du lot -----------------------------------------------------
% Reprises telles quelles par le script de rendu, qui les affiche en tête de
% la vue de lecture. Elles fixent ce que le fichier prétend couvrir : sans
% elles, une transcription flotte sans point d'attache dans un fonds de seize
% mille feuillets.

\newcommand{\folder}[1]{\gdef\grothFolder{#1}}
\newcommand{\batch}[1]{\gdef\grothBatch{#1}}
\newcommand{\leaves}[2]{\gdef\grothFirst{#1}\gdef\grothLast{#2}}
\newcommand{\foldertitle}[1]{\gdef\grothFoldertitle{#1}}
\gdef\grothFolder{?}\gdef\grothBatch{?}\gdef\grothFirst{?}\gdef\grothLast{?}\gdef\grothFoldertitle{}

% --- Le résumé --------------------------------------------------------------
%
% Il ouvre la lecture modernisée, sous le titre « Résumé », et il est composé
% en petit. Ce n'est pas un ornement : c'est le seul passage du document qui
% ne soit pas des mathématiques, et un lecteur doit pouvoir le reconnaître
% avant de l'avoir lu — puis le sauter s'il connaît déjà le sujet, ou s'y
% arrêter s'il ne le connaît pas. Le filet qui le suit marque la reprise du
% corps.

\newenvironment{resume}
  {\small\setlength{\parskip}{0.45em}}
  {\par\medskip\hrule\medskip}

% --- Mots-clés ---------------------------------------------------------------
%
% En anglais, en clôture du résumé de la lecture modernisée : le vocabulaire
% moderne sous lequel chercher ce que les feuillets construisent. Le manifeste
% du site (`npm run manifest`) les extrait pour étiqueter la cote entière —
% cette ligne est la source unique des tags, il n'y a pas de fichier à côté.
\newcommand{\keywords}[1]{\par\smallskip\noindent
  {\footnotesize\textsc{Keywords} --- \textit{#1}}\par}

% --- L'appareil critique ----------------------------------------------------

% Le numéro de feuillet, en marge. C'est l'ancre qui permet de régler un
% désaccord sur une formule en regardant *un* feuillet plutôt qu'une chemise
% de six cents. Le site s'en sert aussi pour tourner les pages du fac-similé
% quand on fait défiler la transcription.
\newcommand{\leaf}[1]{%
  \marginnote{\footnotesize\color{gray!70}\textsf{#1}}%
  \label{leaf:#1}\ignorespaces}

% Illisible : on ne devine pas. Un mot inventé qui se lit comme les autres est
% le pire résultat possible.
\newcommand{\ill}{\textcolor{red!60!black}{[\,\ldots\,]}}

% Lecture douteuse : le mot est donné, et signalé comme incertain.
\newcommand{\uncertain}[1]{\uline{#1}}

% Ajout de l'éditeur — une lettre manquante, un mot sous-entendu.
\newcommand{\add}[1]{\textcolor{blue!50!black}{[#1]}}

% Biffé par Grothendieck lui-même. Ce qu'il a rayé fait partie du document :
% une rature dit souvent où la pensée a changé de direction.
\newcommand{\struck}[1]{\sout{#1}}

% Note du transcripteur, sur l'état matériel du feuillet.
\newcommand{\note}[1]{\par\smallskip{\small\color{gray!60!black}\textit{[#1]}}\par\smallskip}

% Note marginale de Grothendieck, distincte du corps du texte.
\newcommand{\marginal}[1]{\par\smallskip\noindent\hspace{1em}%
  {\small\color{gray!60!black}#1}\par\smallskip}

% --- Titre ------------------------------------------------------------------

\AtBeginDocument{%
  \begin{center}
    {\large\bfseries Fonds Alexandre Grothendieck --- cote n\textsuperscript{o}~\grothFolder}\\[2pt]
    {\itshape\grothFoldertitle}\\[4pt]
    {\small feuillets \grothFirst--\grothLast\ (lot \grothBatch)}\\[2pt]
    {\footnotesize\color{gray!60!black}Transcription établie d'après le fac-similé numérisé
     par l'Université de Montpellier.\\
     Les lectures incertaines sont soulignées, les illisibles notées [\,\ldots\,],
     les ajouts entre crochets.}
  \end{center}
  \vspace{1em}}

\endinput


\folder{11}
\batch{2}
\pages{21}{40}
\dating{[à partir de 1961]}
\watermark{Édition de démonstration}
\foldertitle{Formalisme des correspondances : notes manuscrites (s.d.)}

\begin{document}

\section*{Sous-anneaux bivariants sur $\mathcal{V}$}

\note{Les titres de section reprennent ses propres mots en tête de page. La
page 22 reprend les sous-foncteurs $A' \subset A$ stables de la page 20 (lot
précédent) ; aucune phrase ne franchit la limite des pages 20 et 21.}

\page{22}

\struck{\ill{} Sommes}

Sous-anneaux bivariants sur $\mathcal{V}$. [satisfaisant \ill{} compte les
form.\ de proj.\ et d'échange]

Intersection de tels \uncertain{est} idem ; définition du foncteur « stable »
engendré. Soit $A_{sp} \subset A$ le plus petit ss-foncteur stable

engendré.
\note{La ligne s'arrête après « stable » et « engendré. » est seul sur la
suivante : ce qui engendre $A_{sp}$ n'est pas dit sur la page.}

\emph{NB} $A_{sp}(X)$ contient les $f_{*}(U_{Y})$ pour $f : Y \to X$ dans
$\mathcal{V}$. Dans beaucoup de cas importants, il est $=$ identique à
\uncertain{l'ens.} des $f_{*}(U_{Y})$ pour tous les $f : Y \to X$ dans
$\mathrm{Ob}(\mathcal{V}_{/X})$. Cela \uncertain{marche} p.\ ex.\ si
$\mathcal{V}$ la catégorie des schémas projectifs lisses \ill{} sur $k$ alg.\
clos, dans les cas envisagés plus haut, — dans car.\ 0 si on admet la
résolution des \emph{singularités}.
\note{$U_{Y}$ est l'élément unité de $A(Y)$, noté ailleurs $1_{Y}$.}

\marginal{$A' \subset A$ donné, \emph{stable}}

Une application $A(X_{1}) \times \dots \times A(X_{n}) \to A(X)$ est dite
\emph{du type} $A'$ si elle est de la forme
\[
x_{1}, \dots, x_{n} \longmapsto u(x_{1} \boxtimes \dots \boxtimes x_{n})
\]
avec $u \in A'(X_{1} \times \dots \times X_{n} \times X)$. En
\struck{\ill{}} particulier, \uncertain{cas} \uncertain{défini} si $n = 1$.
\uncertain{Composées} \ill{} \add{de} telle\add{s} applications :
$\varphi$ \ill{} une application $\psi : A(X) \to A(X')$ qui est de type $A'$,
\ill{} est une application qui est du type $A'$. Plus général.
\struck{Plus général. De \ill{}}, \ill{} comp.\ des applications
\struck{des applications $(X_{i})_{i \in I}$}
\[
\begin{aligned}
\varphi(X_{1}, \dots, X_{n}) &: A(X_{1}) \times \dots \times A(X_{n}) \to A(X) \\
\psi(Y_{1}, \dots, Y_{m}) &: A(Y_{1}) \times \dots \times A(Y_{m}) \to A(Y)
\end{aligned}
\]
\ill{} avec $\chi : A(X) \times A(Y) \to A(Z)$

\page{24}

avec $\varphi, \psi, \chi$ du type $A'$, on trouve une application
\struck{$A(X$}
\[
\chi(\varphi, \psi) : A(X_{1}, \dots, X_{n}) \times A(Y_{1}) \times \dots \times A(Y_{m}) \to A(Z)
\]
qui est \struck{$A'$-admissible} du type $A'$ \add{:} de type $A'$.

Dans le produit des deux applications \struck{\uncertain{admissibles}} du type
$A'$ [\ill{}] est \struck{$A'$} du type $A'$, \uncertain{i.e.} $=$ le produit
par un élément de $A'(X)$.

\emph{Application} \struck{Soit $E \subset \mathcal{V}$, et \ill{}} \ill{}
$X \in \mathcal{V}$, supposons donné un $E_{X}$ \uncertain{$\subset$} $A(X)$
(\uncertain{qui peut être vide}).

Alors le plus petit sous-foncteur stable $B \subset A$ contenant $A'$ et les
$E_{X}$ est défini ainsi :
$B(X) =$ ensemble des éléments de \struck{$A(X)$ dans $A'(X)$}
$A'(X) \cup$ \ill{}
de la forme $\varphi(x_{1}, \dots, x_{n})$, où $\varphi$ est une application
$A'$-admissible de $A(X_{1}) \times \dots \times A(X_{n})$, les
$X_{i} \in \mathrm{Ob}\, \mathcal{V}$, $x_{i} \in E(X_{i})$ pour
$i = 1, \dots, n$.

En effet, les $B(X)$ [ainsi définis] sont stables par images directes et
inverses, et contiennent les $1$, et les $A'$ ; réciproquement
\struck{il \ill{}} par \ill{} de $\mathcal{V}$ \ill{} pour \ill{} par
multiplication OK.

\page{26}

\note{Le haut de la page, jusqu'au trait horizontal, est barré de deux longs
traits obliques ; il se lit.}

\begin{tikzcd}
A \times B \arrow[r, "s_{A,B}"] \arrow[d, "f \times g"'] & B \times A \arrow[d, "g \times f"] \\
A' \times B' \arrow[r, "s_{A',B'}"] & B' \times A'
\end{tikzcd}

\[
(g \times f)_{*}\bigl(s(x \otimes y)\bigr) = (g \times f)_{*}(y \otimes x)(-1)^{xy} = g_{*}(y) \otimes f_{*}(x)\,(-1)^{xy + fx}
\]
\[
s\bigl((f \times g)_{*}(x \otimes y)\bigr) = s\bigl(f_{*}(x) \otimes g_{*}(y)\bigr)(-1)^{gx} = g_{*}(y) \otimes f_{*}(x)\,(-1)^{gx + (x+f)(y+g)}
\]
\[
(-1)^{xy + fx + gx + gx + fy + xy + fg} \qquad (-1)^{f(x+y+g)} \qquad (-1)^{fy + gx + fg}
\]
\struck{$= (-1)^{(f+x)(g+y)}$}
\note{Dans la première formule de cette ligne, les termes $xy$, $gx$, $gx$, $xy$
sont barrés un à un, laissant $fx + fy + fg$ ; dans la dernière, un $xy$ est
barré devant $fy$. Les exposants $fx$ et $gx$ des deux lignes précédentes sont
cerclés.}

\section*{A}

\note{Un grand « A » souligné ouvre la section ; « $\mathcal{V}^{\circ}$ » est
écrit au-dessus du « 1) ».}

1) Catégorie $\mathcal{V}^{\circ}$ des alg.\ de Poincaré $A$ sur $k$ : alg.\
graduées \uncertain{anticommutatives} de t.f.\ à degrés $\geq 0$, munie d'un hom.\
$\kappa_{A} : A \to k$ de degré $d^{\circ}(A)$ dépendant de $A$ tel que
$(x, y) \mapsto \kappa_{A}(xy)$ soit une forme strictement non
dégénérée —
\note{Il écrit $k_{A}$ ; la lettre est notée $\kappa_{A}$ pour la distinguer de
l'anneau $k$, comme dans le lot 3. « strictement » est ajouté au-dessus de
« non ».}
\struck{\ill{} de $\mathcal{V}^{\circ}$ \ill{}}

Les hom.\ \add{de} $\mathcal{V}^{\circ}$ : ce sont les hom.\ d'algèbres
graduées, sans plus (donc il n'est pas question dans la définition de
$\mathcal{V}^{\circ}$ \uncertain{hom.}\ des $\kappa_{A}$).

\emph{NB} Un morph.\ de cette catégorie n'est pas nécessairement compatible
avec les $\kappa_{A}$, $\kappa_{B}$.

2) Foncteur contravariant
$\varphi : (\mathcal{V}^{\circ})^{\circ} = \mathcal{V} \to \mathrm{Modgrad}(k)$
a) sur les objets : $\varphi(A) = A$ \uncertain{dénuée de la} multiplication

b) sur les morphismes : soit $f^{*} : B \to A$, on pose
$\varphi(f^{*}) = f_{*} : A \to B$ défini par la condition que
\[ \kappa_{A}\bigl(f^{*}(b)\, a\bigr) = \kappa_{B}\bigl(b\, f_{*}(a)\bigr) \]
On vérifie $\varphi(g^{*} f^{*}) = \varphi(f^{*})\, \varphi(g^{*})$,
$\varphi(\mathrm{id}) = \mathrm{id}$, et formule de proj.
\[ f_{*}(a)\, b = f_{*}\bigl(a\, f^{*} b\bigr) \]
et \struck{$\kappa_{B}(f_{*}(a))$} $\kappa_{B} f_{*} = \kappa_{A}$.

\page{28}

Donc on trouve un foncteur \struck{contravariant}
\[ \mathcal{V} \longrightarrow \text{algèbres bivariantes graduées} \]

3) Dans $\mathcal{V}^{\circ}$ \struck{il y a \ill{}} il y a objet initial :
$k$.

Si $A, B \in \mathcal{V}^{\circ}$, il existe un coproduit dans
$\mathcal{V}^{\circ}$ (produit dans $\mathcal{V}$), savoir le produit
tensoriel gauche $A \otimes B$. Si $f^{*} : A \to A'$, $g^{*} : B \to B'$,
alors $f^{*} \otimes g^{*} : A \otimes B \to A' \otimes B'$ défini par
\[ f^{*} \otimes g^{*}(a \otimes b) = f^{*}(a) \otimes g^{*}(b) \]
et l'hom.\ $A \otimes B$ [\uncertain{correspondant} \ill{}] des coproduits.
\struck{\ill{}} On convient de \ill{} défini par (produit tensoriel gauche
d'homom.\ homogènes de degrés gradués)
\[ \kappa_{A \otimes B} = \kappa_{A} \otimes \kappa_{B} \]
\marginal{produit tensoriel \emph{gauche} d'homom.\ homogènes des produits
gradués}
\marginal{i.e.\ $\kappa_{A \otimes B}(x \otimes y) = \kappa_{A}(x)\,\kappa_{B}(y)\,(-1)^{Bx}$}
\struck{\ill{}}
\[ \varphi(f^{*} \otimes g^{*}) = \varphi(f^{*}) \otimes \varphi(g^{*}) \]
\note{Le passage est surchargé : deux mots encadrés et barrés à gauche de la
formule, des renvois qui se croisent ; l'ordre donné ici est celui des flèches.}

i.e.
\[ (*) \qquad (f \times g)_{*}(a' \otimes b') = f_{*}(a') \otimes g_{*}(b')\,(-1)^{g a'} \]
\emph{Mais c'est \emph{faux}} si on ne suppose pas $f$ et $g$ de degrés pairs.

Les deux membres de formule $(*)$ diffèrent par un facteur
\[ (-1)^{ag + bf + Bf} \]

Si dans la définition de $\kappa_{A \otimes B}$ on n'avait pas mis de signe
$(-1)^{Bx}$, les signes qui \uncertain{viendraient} comme \ill{} seraient
\[ (-1)^{a'g + bf} \]

Si en plus on se \uncertain{dispensant} des \struck{produits} tordus, et qu'on
\uncertain{prend} \ill{} le signe $(-1)^{g a'}$, on trouverait un signe
\[ (-1)^{bf} = (-1)^{B'f} \]

Si on prenant $\kappa$ tordu, $\otimes$ ordinaire, le signe serait
\[ (-1)^{bf + (AB + A'B')} \]

C'est donc le signe \ill{} prend tout ordinaire
($\kappa_{A \otimes B} = \kappa_{A} \otimes \kappa_{B}$ et
$f_{*} \otimes g_{*}$) qui est le plus simple. On trouve donc
\[
\boxed{\ (f \times g)_{*} = \bigl(f_{*} \overset{\mathrm{ord}}{\otimes} g_{*}\bigr)\bigl(\mathrm{id}_{A} \otimes [(-1)^{B}]_{B'}^{\deg f}\bigr)\ }
\]
\note{Le second facteur est surchargé : le crochet $[(-1)^{B}]$, l'indice $B'$
et l'exposant « $\deg f$ » se lisent, leur disposition exacte est
\uncertain{douteuse}.}
\struck{\ill{}} Il n'y a aucun signe dès que $f$ est pair.

\page{30}

Pour la suite, on se borne aux $f = (f^{*}, f_{*})$ de \emph{degré pair},
donc on trouve $(f \times g)_{*} = f_{*} \overset{\mathrm{ord}}{\otimes} g_{*}$.

Considérons l'isom.\ de symétrie de $A \otimes B$ sur $B \otimes A$
[considérés comme produits cart.\ dans $\mathcal{V}$], on vérifie que c'est
$s(x \otimes y) = (-1)^{xy}\, y \otimes x$. On a bien, comme il convient
\[
(g \times f)_{*}\bigl(s(x \otimes y)\bigr) = s\bigl((f \times g)_{*}(x \otimes y)\bigr)
\]
\[
\begin{aligned}
s(x \otimes y) &= y \otimes x\,(-1)^{xy}, & (f \times g)_{*}(x \otimes y) &= f_{*}(x) \otimes g_{*}(y), \\
g_{*}(y) \otimes f_{*}(x)\,(-1)^{xy} &\ \Vert & g_{*}(y) \otimes f_{*}(x)\,&(-1)^{(x+f)(y+g)} .
\end{aligned}
\]
\note{Sur la page, les deux membres sont développés sous des accolades ; les
deux derniers termes sont égaux puisque $f$ et $g$ sont pairs.}

Prouvons les formules de \uncertain{permutation}

\begin{tikzcd}
X \arrow[d, "f"'] & X \times S \arrow[l, "q"'] \arrow[d, "f'"] \\
Y & Y \times S \arrow[l, "p"]
\end{tikzcd}

\[
\left\{
\begin{aligned}
p^{*}\bigl(f_{*}(x)\bigr) &= f'_{*}\bigl(q^{*}(x)\bigr) \\
f^{*}\bigl(p_{*}(m)\bigr) &= q_{*}\bigl(f'^{*}(m)\bigr)
\end{aligned}
\right.
\]
\note{Sous la flèche $p$, une étoile est ajoutée puis surchargée : $p^{*}$.}

\begin{tikzcd}
B \arrow[r, "q^{*}(b) = b \otimes 1"] & B \otimes R \\
A \arrow[u, "f^{*}"] \arrow[r, "p^{*}(a) = a \otimes 1"'] & A \otimes R \arrow[u, "f^{*} \otimes \mathrm{id}_{R} = f'^{*}"']
\end{tikzcd}

\note{« $x \in B$ » est écrit devant le $B$ du carré. Au coin supérieur droit,
un $A$ est surchargé en $R$.}

\[ f_{*}(x) \otimes 1 = (f_{*} \otimes 1)(x \otimes 1) \qquad \text{OK} \]
\[
f^{*}\bigl((\mathrm{id}_{A} \otimes \kappa_{R})(a \otimes r)\bigr) \overset{?}{=} (\mathrm{id}_{B} \otimes \kappa_{R})(f^{*} \otimes \mathrm{id}_{R})(a \otimes r)
\]
\[
\kappa(r)\, a \ \mapsto\ \kappa(r)\, f^{*}(a) \qquad\qquad f^{*}(a)\, r \ \mapsto\ \kappa(r)\, f^{*}(a) \qquad \text{O.K.}
\]
\note{Les deux membres sont calculés sous des accolades : $(\mathrm{id}_{A}
\otimes \kappa_{R})(a \otimes r) = \kappa(r) a$, $(f^{*} \otimes
\mathrm{id}_{R})(a \otimes r) = f^{*}(a) r$ ; sous la première, un premier essai
est biffé.}

\[
X \xrightarrow{\ f\ } Y \xrightarrow{\ g\ } Z \qquad\qquad A \xleftarrow{\ f^{*}\ } B \xleftarrow{\ g^{*}\ } C
\]

\begin{tikzcd}
 & X \arrow[dl, "\Gamma_{gf}"'] \arrow[dr, "\Gamma_{f}"] & \\
X \times Z \arrow[dr, "\Gamma_{f} \times \mathrm{id}_{Z} = \lambda"'] & & X \times Y \arrow[dl, "\mathrm{id}_{X} \times \Gamma_{g} = \lambda'"] \\
 & X \times Y \times Z &
\end{tikzcd}

\begin{tikzcd}
 & A & \\
A \otimes C \arrow[ur] & & A \otimes B \arrow[ul] \\
 & A \otimes B \otimes C \arrow[ul, "\lambda^{*}"] \arrow[ur, "\lambda'^{*}"'] &
\end{tikzcd}

\note{Les éléments sont écrits le long du diagramme : $a \otimes b \otimes c$
sous $A \otimes B \otimes C$, $(a f^{*}(b)) \otimes c$ près de $A \otimes C$,
$a \otimes b\, g^{*}(c)$ près de $A \otimes B$, et au-dessus de $A$
\uncertain{$a f^{*}(b) \cdot f^{*} g^{*}(c)$}. Une flèche biffée monte d'un
symbole raturé vers $A \otimes C$.}

\[ \therefore\ \lambda'^{*}\bigl(\lambda_{*}(u)\bigr) = (\Gamma_{f})_{*}\bigl(\Gamma_{gf}^{*}(u)\bigr) \]
et formule symétrique

\emph{Question} Trouver un cas général commun où la formule d'échange soit
vraie. Suffit-il que le produit fibré soit dans $\mathcal{V}$ et que les dim.\
soient correctes ?

\section*{Catégorie des algèbres bivariantes (sur un anneau $k$)}

\page{33}

\emph{Catégorie des algèbres bivariantes} (sur un anneau $k$) [munie d'une
forme $\kappa_{A} : A \to k$]

C'est la catégorie dont les objets sont les $k$-algèbres, et les morphismes de
$A$ dans $B$ les couples $f = (f_{*}, f^{*})$, où $f^{*} : B \to A$ est un
hom.\ d'algèbres [\ill{}], et où $f_{*} : A \to B$ est un hom.\ $B$-linéaire,
i.e.\ un hom.\ additif \uncertain{satisfaisant} la « formule de projection »
\[
\bigl[\, f_{*}\bigl(x f^{*}(y)\bigr) = f_{*}(x)\, y \,\bigr], \quad f_{*}\bigl(f^{*}(y)\, x\bigr) = y\, f_{*}(x)
\]
\note{Une note entre crochets, reliée à la première formule par une flèche :
« signe $(-1)^{\delta j}$ \ill{} dans le cas gradué ; \uncertain{si la deuxième
formule est vraie}, $y$ de degré $j$, $f_{*}$ de degré $\delta$ ».}
et de plus : $\kappa_{B} f_{*} = \kappa_{A}$.

\marginal{il faut choisir entre l'une et l'autre formule !}

On \ill{} une variante graduée, et on aura à travailler dans diverses
sous-catégories pleines, ou même non pleines.

Dans le cas gradué anticommutatif (où $f^{*}$ est donc hom.\ d'alg.\ graduées
[\emph{en degré pair}], et où $f_{*}$ gradué) \ill{} (les \ill{}), les deux
formules $(*)$ sont équivalentes.

\marginal{\emph{Remarque} Pour que $(f_{*}, f^{*})$ soit un isomorphisme, il faut
et il suffit que $f^{*}$ soit \struck{\uncertain{surjectif}} \uncertain{bijectif}
\ill{} $f_{*}$ \ill{}, i.e.\ $f^{*}$ un isom.\ et $f_{*}(1)$ inversible ; les
« bons » isom.\ \uncertain{sont} ceux où $f_{*}(1) = 1$, i.e.\ $f^{*}$ et
$f_{*}$ inverses l'un de l'autre.}
\note{La remarque est écrite en oblique dans la marge gauche, dans une bulle
reliée au paragraphe précédent.}

Soit $A$ une algèbre munie d'une forme
\[ \kappa_{A} : A \to k \]
\struck{que} dans la forme bilinéaire
\[ (x, y) \longmapsto \kappa(xy) \]
(= \uncertain{deux} !),
sur $A \times A$ \struck{soit séparante, dans} définissant un hom.\
\struck{injectif}
\[
\varphi_{A} : A \to \check{A} \qquad \bigl(\check{A} = \mathrm{Hom}_{k\text{-mod}}(A, k)\bigr), \qquad \langle y, \varphi_{A}(x) \rangle = \kappa(yx) .
\]
Supposons que $\varphi_{B}$ soit injectif, $B$ \ill{} de même muni de
$\kappa_{B} : B \to k$, et soit $f^{*} : B \to A$ tel que l'hom.\ transposé
${}^{t}f^{*} : \check{A} \to \check{B}$ applique $\mathrm{Im}\, \varphi_{A}$ dans
$\mathrm{Im}\, \varphi_{B}$ [ce qui n'est pas une restriction si
$\varphi_{B} : B \xrightarrow{\sim} \check{B}$ !]. Alors l'application induite,
[unique], $f_{*}$ est donc caractérisée par
\[ (***) \qquad \kappa\bigl(y\, f_{*}(x)\bigr) = \kappa\bigl(f^{*}(y)\, x\bigr) \]
et satisfait de ce fait la 2ème relation $(*)$, [dans le
premier cas commutatif, ou gradué anticommutatif] [De plus, il satisfait
$(**)$ (faire $y = 1$ dans $(***)$)] quand $f^{*}$ est hom.\ d'alg.\ unitaire
\ill{} près]. \struck{Supposons \ill{} $A$, $B$} Réciproquement, \ill{} hom.\
$(f^{*}, f_{*})$ de $(A, \kappa_{A})$ dans $(B, \kappa_{B})$, $f_{*}$
satisfait $(***)$, \uncertain{complètement} déterminé par $f^{*}$.
\note{« injectif » est lu en surcharge au-dessus d'un mot barré, dans
« Supposons que $\varphi_{B}$ soit injectif ».}

\page{35}

gradué anticommutatif, $f^{*}$ homogène. Supposons qu'il existe
\struck{deux} entier $d = d(A)$ tel que $\kappa_{A}$ soit composé
\[ A \to A^{d} \xrightarrow{\ \bar{\kappa}_{A}^{(d)}\ } k , \]
et de même $d' = d(B)$. Alors si $x$ est de degré $i$, la forme
$\varphi_{A}(x)$ est nulle en degré $\neq d - i$ ; donc
$\langle f^{*}(y), \varphi_{A}(x) \rangle = \kappa\bigl(f^{*}(y), x\bigr)$ est
nulle \struck{\ill{}} en degré $\neq d - i$ ; \struck{donc si
$x$ \ill{}}, si $x'$ \struck{\ill{}} est hom.\ de degré $i'$, la forme
$\varphi_{B}(x')$ est nulle en degré $\neq d' - i'$, donc on doit avoir (si
$\varphi_{B}(x') \neq 0$)
\[ d - i = d' - i' \qquad \text{i.e.} \qquad i' = i + (d' - d) . \]
Si donc on pose $f_{*}(x) = \sum_{j} x'_{j}$, on voit que l'effet de la forme
$\varphi_{B}(f_{*}(x))$ sur $B^{d'-j}$ est égal à $\varphi_{B}(x'_{j})$, donc est
nul si \struck{$j$} $d' - j \neq d - i$, donc dans ce cas $x'_{j} = 0$. Donc
$f_{*}$ \emph{augmente les degrés de $d' - d$}.

Nous travaillons dans la catégorie [bivariante] des algèbres graduées
anticommutatives, munies d'un degré $d$, les $f_{*} : A \to B$ étant
assujetties à [augmenter] des degrés $d(B) - d(A)$ [et $\kappa_{A}$ se
factorisant par $A^{d(A)}$]. \uncertain{Les sous-catégories} pleines des $A$
tels que $d(A)$ pair et commutatives [et il y a \ill{} éventuellement de se
borner aux algèbres $A$ qui sont à degrés \struck{\ill{}} dans $[0, d(A)]$]. Il
n'est pas question de se borner aux « algèbres de Poincaré », même en
\uncertain{homologie modulo torsion} ; \ill{} les formes sur $A$, car les
constructions universelles qui se font \ill{} de \ill{} \ill{}
\uncertain{fassent sentir} !

\page{37}

\struck{Soient $A$ et $B$ deux algèbres}

Le produit tensoriel gradué gauche $A \otimes B$ est un foncteur sur la
catégorie $C$ [bivariante] des algèbres graduées. Le produit tensoriel
(malheureusement !) ne s'interprète pas comme un produit dans la catégorie
bivariante : on n'a pas \ill{} d'hom.\ de la catégorie
$f : A \to A \otimes A$, $f = (f^{*}, f_{*})$, qui \uncertain{jouerait} le rôle
d'une application diagonale. On trouve néanmoins que c'est un produit (dans
$C$) quand $A$, $B$ sont toutes deux des alg.\ de Poincaré ….

\emph{Algèbres bivariantes simpliciales} : c'est un objet simplicial (pas
$\frac{1}{2}$ simplicial !) de la catégorie des algèbres bivariantes, donc un
foncteur contravariant de la catégorie des ens.\ finis non vides dans $C$. Il
est dit \emph{spécial} « s'il transforme sommes en produit tensoriel », i.e.\
si pour $I = I' \sqcup I''$, les hom.
\[ A(I') \xrightarrow{\ \alpha^{*}\ } A(I), \qquad A(I'') \xrightarrow{\ \beta^{*}\ } A(I) \]
[où $\alpha : I' \to I$, $\beta : I'' \to I$] définissent un
\emph{isom.}
\[ A(I') \otimes A(I'') \xrightarrow{\ \sim\ } A(I' \sqcup I'') . \]

Si \struck{$A(\Delta^{\circ})$ est une algèbre} on se borne à la catégorie
pleine des Alg.\ biv.\ simpliciales $A(\cdot)$ telles que
$A(\Delta(0)) = A_{0}$ soit une algèbre de Poincaré, on trouve une
\emph{équivalence} de la catégorie des alg.\ biv.\ simpliciales et des alg.\ de
Poincaré.
\note{L'argument écrit $A(\Delta(0))$ se lit mal ; c'est l'algèbre de l'ensemble
à un élément, que la ligne nomme $A_{0}$.}

\page{39}

\emph{Proposition} Pour tout $n \geq 1$, soit $A_{n} \subset H^{*}(X^{n})$ une
ss-\uncertain{algèbre} \struck{\ill{}} stable \emph{par gpe sym.} Conditions
équivalentes :

(i) a) Pour tout \struck{$u \in H^{*}$} $u \in A_{n+m} \subset H^{*}(X^{n} \times X^{m})$
et $x \in A_{n} \in H^{*}(X^{n})$, on a $u(x) \in H^{*}(X^{m})$.

b) Si $x \in A_{n}$, $y \in A_{m}$, alors $x \boxtimes y \in A_{n+m}$ ; pour
tout $n$, $1 \in A_{n}$.

c) Pour \struck{toute} \struck{surjection}
$[1, n+1] \twoheadrightarrow [1, n]$, son graphe dans
$H^{*}(X^{n} \times X^{n+1})$ est $\in A_{2n+1}$.
[Et de même pour les injections $[1, n] \hookrightarrow [1, n+1]$.]

(ii) Les \struck{systèmes des} $A_{n}$ sont des \emph{sous-algèbres}, stables
par images inverses [et images directes] par \emph{projections}
$X^{n+1} \to X^{n}$, et \struck{par images inverses} par images directes par
injections $X^{n} \hookrightarrow X^{n+1}$.

Ces conditions impliquent que les $A_{n}$ sont stables par images inverses et
images directes par toute application simpliciale [$X^{n} \to X^{m}$], et que
\struck{$A_{n}$} \ill{} pour telles applications,
$H^{*}(X^{n} \times X^{m})$ \uncertain{contient} la classe du graphe…

(iii) Considérons $A_{n+m} \subset H^{*}(X^{n} \times X^{m})$ comme un ensemble
d'homomorphismes de $H^{*}(X^{n})$ dans $H^{*}(X^{m})$, on veut que ces
ensembles soient stables par composition ($n, m \geq 0$), contiennent les
identités, soient stables par $\blacksquare$, et contiennent les homomorphismes
images inverses et directes par une application simpliciale.
\note{Le symbole après « stables par » est un carré plein, sans doute le
produit externe $\boxtimes$ de (i) b).}

\end{document}
